\subsection{R-26: a uniform linear homological defect from Jacobi continuants}

\paragraph{Statement R-26.}
For every $m\geq17$, some ternary Jacobi pencil and endpoint vectors
in the notation of R-25 satisfy
\begin{equation}\label{R26:criterion}
 \dim\mathcal C_A(u,v)\geq8m-4.
\end{equation}
They can be chosen with
\[
 A_{ii}=\alpha_i x,\quad A_{i,i+1}=y,\quad
 A_{i+1,i}=\delta_i z,\qquad
 u=(x,0,\ldots,0,y)^t,\quad v=(x,0,\ldots,0,z)^t,
\]
where all $\alpha_i,\delta_i$ are positive integers at most $8m^2$.
No efficiency claim is made for selecting them from this finite grid.

Consequently, for every $m\geq18$,
\begin{equation}\label{R26:linear_defect}
 0\leq\epsilon_m\leq6(m-18),\qquad
 4m^2-8m+110\leq c_m\leq4m^2-2m+2.
\end{equation}
In particular every integer pair
\begin{equation}\label{R26:padding}
 m\geq18,\qquad m<n,\qquad n^2\leq2(m-1)^2+53
\end{equation}
satisfies the strict four-variable dimension comparison of R-14.
This includes every integer $m<n\leq\sqrt2(m-1)$.
The explicit stretched-direction conclusions of G-17/G-18 therefore
apply to all these pairs. The uniform conclusion is a linear bound
on the defect, not its vanishing at every size.

\begin{proof}
We first prove \eqref{R26:criterion}. Write $L=m-1$ and $h=yz$.
For an interval of vertices define its continuant
\[
 P_{[a,b]}(x,h)=\operatorname{per}(A_{[a,b],[a,b]}),\qquad
 P_{[a,a-1]}=1.
\]
Its recurrence is
\begin{equation}\label{R26:continuant}
 P_{[a,b]}=\alpha_bxP_{[a,b-1]}
                   +\delta_{b-1}hP_{[a,b-2]}\quad(b>a),
 \qquad P_{[a,a]}=\alpha_ax.
\end{equation}
The two summands distinguish whether the last vertex is fixed or
paired with its predecessor. An interval of $s$ vertices is
homogeneous of weighted degree $s$ for $\deg x=1$, $\deg h=2$.
Put
\[
 D_i=\prod_{e=i}^{L-1}\delta_e,\qquad
 E_i=\prod_{e=0}^{i-1}\delta_e,
\]
with empty products one.

\paragraph{1. Exact endpoint formulas and character spaces.}
The complementary matching for a tridiagonal matrix consists of
the forced path between the removed row and column, and independent
matchings on the intervals outside that path. Thus the first
auxiliary derivatives of R-25 are exactly
\begin{align}
 a_i&=x y^iP_{[i+1,L]}+D_i z^{m-i}P_{[0,i-1]},\label{R26:a}\\
 b_i&=x E_i z^iP_{[i+1,L]}+y^{m-i}P_{[0,i-1]}.
                                                        \label{R26:b}
\end{align}
These formulas include the endpoints using empty intervals.
For $0<i<L$, differentiating the quadratic coefficient of
$\operatorname{per}(A+tuv^t)$ gives
\begin{align}
 t_i&=2x^2D_i z^{m-i}P_{[1,i-1]}
           +2xy^ihP_{[i+1,L-1]},\label{R26:t}\\
 t'_i&=2x^2y^{m-i}P_{[1,i-1]}
           +2xE_i z^ihP_{[i+1,L-1]}.
                                                        \label{R26:tp}
\end{align}
Indeed, the two selected columns are the endpoints. The other
selected row is either endpoint $0$ or endpoint $L$; the intervening
path is forced, leaving the indicated residual interval. The two
assignments of the rank-one entries give the factor $2$.
At the endpoints, with $M=P_{[1,L-1]}$,
\begin{equation}\label{R26:endpoints}
 t_0=t'_0=2xhM,\qquad t_L=2x^2zM,\qquad t'_L=2x^2yM.
\end{equation}

Let a primitive $m$th root $\omega$ act by
$(x,y,z)\mapsto(x,\omega y,\omega^{-1}z)$.
Equations \eqref{R26:a}--\eqref{R26:tp} show that $a_i,t_i$
have character $i$, and $b_i,t'_i$ have character $-i$ modulo $m$.
The polynomial $P=P_{[0,L]}$ has character zero.
Consequently different character spaces can be added independently.

\paragraph{2. Eight columns in each ordinary character.}
Fix $2\leq r\leq m-2$ and put $j=m-r$.
We will prove that the following eight columns are generically
independent:
\begin{equation}\label{R26:eight_columns}
 xa_r, ya_{r-1}, za_{r+1}, xb_j, yb_{j+1}, zb_{j-1}, t_r, t'_j.
\end{equation}
For this assertion specialize $\alpha_i=1$ and retain the
$\delta_i$ as parameters. Each column has the unique form
$y^r p(x,h)+z^jD_rq(x,h)$, where $p$ has weighted degree $j+1$
and $q$ has weighted degree $r+1$. These two summands have different
integer weights, so they occupy disjoint monomial coordinates.
When the edges are nonzero, dividing the second component by
$D_r$ preserves rank. Write
\begin{equation}\label{R26:K}
 K=\frac{E_j}{D_r}
  =\frac{\prod_{e=0}^{\min(r,j)-1}\delta_e}
         {\prod_{e=\max(r,j)}^{m-2}\delta_e}.
\end{equation}
The eight resulting pairs $(p,q)$, in the order
\eqref{R26:eight_columns}, are
\begin{equation}\label{R26:pair_matrix}
\begin{array}{c|l|l}
 &p&q\\ \hline
1&x^2P_{[r+1,L]}&xP_{[0,r-1]}\\
2&xP_{[r,L]}&\delta_{r-1}hP_{[0,r-2]}\\
3&xhP_{[r+2,L]}&\delta_r^{-1}P_{[0,r]}\\
4&xP_{[0,j-1]}&Kx^2P_{[j+1,L]}\\
5&P_{[0,j]}&K\delta_jxhP_{[j+2,L]}\\
6&hP_{[0,j-2]}&K\delta_{j-1}^{-1}xP_{[j,L]}\\
7&2xhP_{[r+1,L-1]}&2x^2P_{[1,r-1]}\\
8&2x^2P_{[1,j-1]}&2KxhP_{[j+1,L-1]}.
\end{array}
\end{equation}
All entries in this display follow directly from
\eqref{R26:a}--\eqref{R26:tp}; in particular the long path
products have disappeared except for the single scalar $K$.

\paragraph{3. Interpreting cut limits correctly.}
For the moment regard $K$ as an independent nonzero scalar in
\eqref{R26:pair_matrix}. We allow some edges to tend to zero,
while keeping $\delta_r$ and $\delta_{j-1}$ nonzero. Such a limit
is a valid rank witness whenever it is approached by nonzero edges
with \eqref{R26:K} tending to the specified $K$.

There are two situations used below. If all zero edges are between
$\min(r,j)$ and $\max(r,j)-1$, they do not occur in the reduced
ratio \eqref{R26:K}; its value is the ratio of the fixed outer
products. If there are additionally equally many zero edges in
the left outer interval $[0,\min(r,j)-1]$ and the right outer
interval $[\max(r,j),m-2]$, with at least one on each side, then
any prescribed nonzero $K$ is possible. To see this, replace each
zero edge by $\tau$ times a nonzero constant. The powers of
$\tau$ in the outer products cancel. Adjust the constant of one
left edge to give the desired ratio; the other constants can be
one. For every $\tau\ne0$ all edges are nonzero, and the entries
of \eqref{R26:pair_matrix} have the stated finite limit.

Thus a nonzero limit minor proves a nonzero minor before taking
the limit. This is a rational row normalization followed by a
limit; it is not a claim about the rank of the literal pencil
obtained by simply setting its edges to zero.

\paragraph{4. Two exact insertion operations.}
The following operations on the normalized matrix preserve its
rank. They use only \eqref{R26:continuant}: inserting an isolated
vertex with diagonal $x$ at a zero edge multiplies every interval
continuant containing that vertex by $x$.

First suppose $r,j\geq6$, with zero edges at positions $2$ and
$m-4$. Insert one isolated vertex at each of these two cuts,
replacing each zero edge by two zero edges. Then
\[
 (m,r,j)\longmapsto(m+2,r+1,j+1).
\]
Every interval in the $p$ column of \eqref{R26:pair_matrix}
contains exactly one new vertex: prefixes contain the left one
and suffixes the right one. The same is true in the $q$ column.
The cuts are away from every endpoint shifted by one in that
display, since $\min(r,j)\geq6$. Accordingly every pair becomes
$(xp,xq)$. The distinguished nonzero edges $\delta_r$,
$\delta_{j-1}$ and the adjacent coefficients are merely
reindexed. The new outer zero-edge counts remain equal, so $K$
can retain its previous value by the preceding limit argument.
Multiplication by $x$ is injective in both components and
preserves the rank. The new pattern still has a left cut at $2$
and a right cut at the new position $(m+2)-4$, allowing repetition.

Second suppose $j-r\geq4$, with a zero edge at $r+1$.
Insert an isolated vertex there. This sends $(m,r,j)$ to
$(m+1,r,j+1)$. Every $p$ interval in
\eqref{R26:pair_matrix} contains the new vertex, including
$[r+2,L]$ and $[0,j-2]$; every $q$ interval avoids it.
Thus the pairs become $(xp,q)$. The cut lies in the middle
interval absent from the reduced ratio \eqref{R26:K}, so $K$
is unchanged. If instead $r-j\geq4$, a cut at $j+1$ gives
$(m+1,r+1,j)$ and pairs $(p,xq)$. The same endpoint check applies.
Both operations preserve injectivity and leave a cut at the
appropriate position for repetition.

\paragraph{5. Complete finite core certificate.}
It remains to supply finitely many starting matrices. Work
modulo $101$. Before imposing cuts, let $\delta_i$ be the
residue of $7i^2+13i+36$ in $\{1,\ldots,100\}$, replacing a
zero residue by $1$. Take the following cores.

For the nine pairs $(r,j)=(6+\max(-g,0),6+\max(g,0))$,
$-4\leq g\leq4$, put $\delta_2=\delta_{m-4}=0$ and $K=14$.
If $|g|=4$, also put $\delta_{\min(r,j)+1}=0$.
For the eight pairs $r+j=17$, $\min(r,j)\in\{2,3,4,5\}$,
put only $\delta_{\min(r,j)+1}=0$ and define $K$ by the exact
rational outer-product ratio \eqref{R26:K}. The residues of
these latter $K$ values are shown below. The edges appearing
in denominators in \eqref{R26:pair_matrix} are nonzero in all
these cases.

Expand the positive components in the ordered basis
$x^{j+1-2k}h^k$, $0\leq k\leq\lfloor(j+1)/2\rfloor$, followed
by the negative components in $x^{r+1-2k}h^k$,
$0\leq k\leq\lfloor(r+1)/2\rfloor$. Number these coefficient
rows from zero. All eight columns are used, in the displayed
order. The following ordered row lists have the stated nonzero
determinants modulo $101$:
\begin{center}
\begin{tabular}{rrrlr}
$r$&$j$&$K\bmod101$&ordered rows&determinant\\ \hline
10&6&14&$(0,1,4,2,3,5,6,7)$&73\\
9&6&14&$(0,1,4,2,3,5,6,7)$&58\\
8&6&14&$(0,1,4,2,3,5,6,7)$&16\\
7&6&14&$(0,1,4,2,3,5,6,7)$&8\\
6&6&14&$(0,1,4,2,3,5,6,7)$&53\\
6&7&14&$(0,1,5,3,2,6,7,8)$&51\\
6&8&14&$(0,1,5,3,4,6,2,7)$&72\\
6&9&14&$(0,1,6,3,4,7,2,8)$&20\\
6&10&14&$(0,1,6,3,4,7,2,8)$&90\\
2&15&34&$(0,1,9,3,4,10,6,7)$&51\\
15&2&34&$(0,1,2,3,4,5,6,7)$&16\\
3&14&61&$(0,1,8,3,4,9,6,7)$&88\\
14&3&61&$(0,1,3,4,2,5,6,7)$&93\\
4&13&88&$(0,1,8,3,4,9,6,7)$&23\\
13&4&88&$(0,1,3,2,4,5,6,7)$&24\\
5&12&82&$(0,1,7,3,4,8,6,2)$&83\\
12&5&82&$(0,1,4,2,3,5,6,7)$&40
\end{tabular}
\end{center}
This is a finite certificate specified entirely by the recurrence
\eqref{R26:continuant}, the pair matrix
\eqref{R26:pair_matrix}, the integer parameter rule, and the
row lists. The source \texttt{check\_jacobi\_cores.py} and its
manifest give every coefficient matrix. Integer Bareiss
determinants of the listed matrices were also separately
computed and give exactly these residues. An independent agent
then reconstructed all seventeen coefficient matrices by direct
enumeration of path matchings, rather than the continuant recurrence,
and recomputed every selected determinant by modular fraction-free Bareiss
elimination. All matrices and residues agreed. Each nonzero residue
is a nonzero rational minor in characteristic zero.

For any target $r,j\geq6$, first choose the core with the same
signed gap, truncated to $[-4,4]$. Repeated paired insertions
raise the smaller index from $6$ to $\min(r,j)$, preserving
the gap. If the desired absolute gap exceeds $4$, repeated
central insertions then raise the larger index to its target.
For $\min(r,j)\in\{2,3,4,5\}$ and $r+j\geq17$, start at
the corresponding size17 core and use central insertions.
Every ordinary character $2\leq r\leq m-2$, $m\geq17$, is
therefore covered. Its eight-column minor is nonzero for some
nonzero edge parameters, so it is nonzero on a nonempty Zariski
open set of the full Jacobi parameter space.

\paragraph{6. The three exceptional characters.}
Here put all $\delta_i=1$ and all $\alpha_i=1$ except
$\alpha_{m-2}=2$. Let $C_0=1,C_1=x$ and
$C_s=xC_{s-1}+hC_{s-2}$ denote the unperturbed continuants.
Consecutive $C_s,C_{s-1}$ are coprime in $\mathbb C[x,h]$.
Indeed, a common irreducible divisor divides $hC_{s-2}$.
If it does not divide $h$, repeat the recurrence until it divides
$C_0=1$, a contradiction. If it divides $h$, it is associated
to $h$, whereas $C_s(x,0)=x^s$ shows that $h$ is not a factor.
For the degree arguments below, $x,h$ have weights $1,2$.

Set
\[
 S_0=P_{[1,L]},\quad S_1=P_{[2,L]},\quad
 F=P_{[0,L-1]},\quad M=P_{[1,L-1]},\quad
 A_0=C_{m-3},\quad B_0=C_{m-4}.
\]
The single changed diagonal gives
\begin{align*}
 S_0&=(2x^2+h)A_0+xhB_0,&
 S_1&=xA_0+(h+x^2)B_0,\\
 F&=(2x^2+h)A_0+2xhB_0,&
 M&=2xA_0+hB_0.
\end{align*}

In character zero, the available columns span
\begin{equation}\label{R26:zero_basis}
 x^2S_0,\quad xhS_1,\quad hF,\quad xhM,\quad xy^m,\quad xz^m.
\end{equation}
For clarity, $xP=x^2S_0+xhS_1$;
$xa_0=x^2S_0+xz^m$ and $za_1=xhS_1+xz^m$;
$xb_0=x^2S_0+xy^m$; $ya_L=xy^m+hF$;
and $t_0=2xhM$. These identities recover all six displayed
forms, since $2\ne0$.
The two extreme monomials are independent of the balanced
polynomials in $x,h$. To prove that the four balanced forms
are independent, suppose their coefficients are $a,b,c,d$.
Their relation becomes
\begin{align*}
0={}&A_0\big(2ax^4+(a+b+2c+2d)x^2h+ch^2\big)\\
 &+B_0xh\big((a+b)x^2+(b+2c+d)h\big).
\end{align*}
Coprimeness and $\deg A_0=m-3>5$ imply that both coefficient
polynomials vanish. Their coefficients then give
$a=b=c=d=0$. Thus character zero has dimension at least six.

In character one, the space contains the four pure weight-one
forms
\begin{equation}\label{R26:one_basis}
 yP,\qquad x^2yM,\qquad x^2yS_1,\qquad xyF.
\end{equation}
The first two are $yP$ and $t'_L/2$. The identity
$za_2-xa_1-ya_0=-2x^2yS_1$ gives the third, and
$xa_1-xb_L=x^2yS_1-xyF$ gives the fourth.
After removing the factor $y$, a relation with coefficients
$a,b,c,d$ becomes
\begin{align*}
0={}&A_0x\big((2a+2b+c+2d)x^2+(2a+d)h\big)\\
 &+B_0\big(ah^2+(2a+b+c+2d)x^2h+cx^4\big).
\end{align*}
Now $\deg A_0>4$ and coprimeness again force all four
coefficients to vanish. Modulo pure integer weight one, the
columns $xa_1,ya_0,yb_0$ have the independent images
$x^2z^{m-1},yz^m,y^{m+1}$. Character one therefore has dimension
at least seven. Transposition together with $y\leftrightarrow z$
preserves this Jacobi pencil and exchanges $u,v$, so character
$-1$ has the same lower bound. All degree inequalities hold
throughout the stated range $m\geq17$.

\paragraph{7. One common pencil and the uniform consequences.}
Each ordinary-character minor and each exceptional-character
minor is nonzero at some point of the irreducible Jacobi parameter
space with nonzero $\alpha_i,\delta_i$. Their nonempty open
sets have nonempty finite intersection. At one common pencil,
the character dimensions therefore add to
\[
 8(m-3)+6+7+7=8m-4,
\]
proving \eqref{R26:criterion}. This intersection step is essential:
the separate cut specializations are witnesses that the different
minor polynomials are nonzero, not simultaneous zero-edge choices.

There is also a bounded integer choice. Every coefficient entry
of the original columns has parameter degree at most $m$.
Choose one nonzero minor of size eight for each ordinary
character, and minors of sizes six, seven, seven for the other
three. Their product has degree at most
$8m(m-3)+20m=8m^2-4m$. The finite-grid lemma used in R-22
therefore gives a nonzero value with all $2m-1$ parameters in
$\{1,\ldots,8m^2\}$. This proves the asserted bounded choice
without a polynomial-time selection claim.

The exact cancellation and open-intersection proof of R-25 works
with any number $s$ of independent ternary normal columns:
their restrictions are independent of $\ell T_A$, so the same
limit minor gives $c_{m+1}\geq c_m+s$. Applying
\eqref{R26:criterion} yields
\[
 c_{m+1}\geq c_m+8m-4\qquad(m\geq17).
\]
Since the torus upper bound increases by $8m+2$, R-23 gives
$\epsilon_{m+1}\leq\epsilon_m+6$.
Starting from $\epsilon_{18}=0$ in R-25 proves
\eqref{R26:linear_defect}. Finally
\[
 c_m\geq4(m-1)^2+106,
\]
whereas \eqref{R26:padding} implies
$2n^2\leq4(m-1)^2+106<c_m+1$. Thus $c_m+3>2n^2+2$,
which is exactly the strict dimension hypothesis used in R-14
and G-17/G-18. This proves \eqref{R26:padding}.
\end{proof}


\paragraph{Corollary R-26: the eventual rectangular padding range from size14.}
For every $m\geq14$ and every integer $m<n<4m/3$, the padded
permanent and determinant have the explicit rectangular
multiplicity gaps supplied by G-11/G-17, with the eventual-stretch
conclusion of G-18. For every $n\geq4m/3$, G-12 excludes
rectangular gaps. The threshold $14$ is sufficient here; no
optimality assertion is made.

\begin{proof}
For $m\geq18$,
\[
 2(m-1)^2+53-\frac{16m^2}{9}
       =\frac29(m-9)^2+37>0.
\]
Thus every $m<n<4m/3$ lies inside the dimension range
\eqref{R26:padding}. For the four smaller sizes use R-13/R-24:
\begin{center}
\begin{tabular}{rrrr}
$m$&largest integer $n<4m/3$&$c_m$&$2n^2-1$\\ \hline
14&18&680&647\\
15&19&816&721\\
16&21&969&881\\
17&22&1124&967
\end{tabular}
\end{center}
In each case $c_m>2n^2-1$, equivalently
$c_m+3>2n^2+2$. The same strict inequality holds for every
smaller eligible $n$. G-11/G-17 therefore apply in the entire
strict rectangular range, and G-18 gives the stated eventual
stretches. G-12 supplies the exclusion at and above $4m/3$.
\end{proof}
